\documentclass[11pt,a4paper,oneside]{amsbook}
\usepackage[T1]{fontenc}
\usepackage[utf8]{inputenc}
\usepackage{lmodern}
\usepackage[american]{babel}
\usepackage{microtype}
\usepackage{geometry}
\geometry{margin=1.1in}
\makeatletter
\def\input@path{{../}}
\makeatother
\usepackage{sga2-en}
\usepackage[hidelinks]{hyperref}

\addto\captionsamerican{\renewcommand{\chaptername}{Expos\'e}}

\title[SGA~2, Exposé XII (English draft)]{Applications to projective algebraic schemes\\
{\normalsize SGA~2, Exposé XII\\English draft}}
\author{A.\ Grothendieck}
\date{}
\hypersetup{
  pdftitle={SGA 2, Exposé XII (English draft)},
  pdfauthor={Alexander Grothendieck},
  pdfsubject={Unofficial English translation of SGA 2, Exposé XII}
}

\begin{document}
\frontmatter
\maketitle

\chapter*{About this draft}
This is an unofficial English translation of Expos\'e~XII, ``Applications to projective algebraic schemes'',
from SGA~2. The source volume is A.\ Grothendieck
(notes by a group of auditors), with an expos\'e by M.\ Raynaud,
\textit{Cohomologie locale des faisceaux coh\'erents et th\'eor\`emes
de Lefschetz locaux et globaux} (SGA~2), S\'eminaire de g\'eom\'etrie
alg\'ebrique du Bois Marie, 1962. It follows the slightly corrected
SMF recomposition,
\href{https://arxiv.org/abs/math/0511279}{arXiv:math/0511279}.

This draft contains the assigned body of the source, including proofs,
original footnotes, and editor notes (N.D.E.). Statement numbering and
mathematical notation follow the source. References to material outside
this draft retain their original numbers. Apparent mathematical
misprints in the source have been retained; they are recorded in the
accompanying README, separately from the translated text. Scholarly
proofreading remains outstanding.

The translator's contribution is licensed under
\href{https://creativecommons.org/licenses/by-sa/4.0/}{CC BY-SA 4.0}.
The French original remains copyright of its authors and publishers.

\tableofcontents
\mainmatter
\setcounter{chapter}{11}
\chapter[Applications to projective algebraic schemes]
{Applications\\ to projective algebraic schemes} \label{XII}

\renewcommand{\theequation}{\arabic{equation}}

\section[Projective duality theorem and finiteness theorem]
{Projective duality theorem and finiteness theorem\protect\sfootnotemark}\label{XII.1}

The following theorem,\sfootnotetext{The present expos\'e, written in January~1963, is distinctly more detailed than was the oral expos\'e, in June~1962.} essentially contained in FAC\sfootnote{J.-P.~Serre, {\og Faisceaux alg\'ebriques coh\'erents\fg}, \emph{Ann. of Math.} \textbf{61} (1955), p\ptbl197-278.} (except that at the time, Serre did not have at his disposal the language of Ext of sheaves of modules\nde{the reader fond of the History of Mathematics will consult with interest the letter that Grothendieck wrote to Serre on 15 December 1955 and the latter's reply of 22 December of the same year; see \emph{Correspondance Grothendieck-Serre}, edited by Pierre Colmez and Jean-Pierre Serre, Documents Math\'ematiques, vol.~2, Soci\'et\'e Math\'ematique de France, Paris, 2001.}), is the global analogue of the local duality theorem (Expos\'e \Ref{IV}), which was modelled on it.

\begin{theorem} \label{XII.1.1}
Let $k$ be a field, $X=\PP_k^r$ the projective space of dimension $r$ over $k$, $F$~a variable coherent module on $X$; then one has an isomorphism of $\partial$-functors in~$F$:
\steco
\begin{equation} \label{eq:XII.1} {}\H^i(X, F)' \isomfrom \Ext^{r-i}(X;F, \omx^r),
\end{equation}
where one sets
\steco
\begin{equation} \label{eq:XII.2} {}\Omega_{X/k}^r= \Oo_{\PP_k^r}(-r-1).
\end{equation}
\end{theorem}

\begin{remarkstar}
Of course, this Module is also the Module of relative differentials of degree $r$ of $X$ over $k$. In this form, the theorem remains true if $X$ is a proper and smooth scheme over $k$ (for the projective case, see A. Grothendieck, \og \emph{Th\'eor\`emes de dualit\'e pour les faisceaux alg\'ebriques coh\'erents}\fg, S\'eminaire Bourbaki May 1957)\sfootnote{For a more general duality theorem, \cf the Hartshorne Seminar cited at the end of \Exp \Ref{IV}.}. When $F$ is locally free, one recovers Serre's duality theorem $\H^i(X, F)' \isomto \H^{r-i}(\SheafHom_{\OX}(F, \omx^r))$. Theorem \Ref{XII.1.1} (which moreover recovers the case of an $X$ projective over $k$, as in {\loccit}) will suffice for our purposes.
\end{remarkstar} The homomorphism (\Ref{eq:XII.1}) is deduced from the Yoneda pairing
\steco
\begin{equation} \label{eq:XII.3} {}\H^i(X, F) \times \Ext^{r-i}(X;F, \omx^r) \to \H^r(X, \omx^r),
\end{equation}
and from a well-known isomorphism (\cf FAC, or \EGA III 2.1.12):
\steco
\begin{equation} \label{eq:XII.4} {}\H^r(X, \omx^r) = \H^r(\PP_k^r, \Oo_{\PP_k^r}(-r-1)) \isomto k.
\end{equation}
To show that (\Ref{eq:XII.1}) is an isomorphism, one proceeds as in the case of the local duality theorem, noting that $\H^r(X, F)$ as a functor in $F$ is right exact (since $\H^n(X, F)=0$ for $n>r$), and that every coherent Module is isomorphic to a quotient of a direct sum of Modules of the form $\Oo(-m)$, with $m$ large. This reduces us by descending induction on $i$ to carrying out the verification for a sheaf of the form $\Oo(-m)$, where this is contained in the well-known explicit computations (FAC, or \EGA III 2.1.12). One may moreover assume $-m \leq -r-1$, in which case $\H^i(X, \Oo(-m))=0$ for $i\neq r$.

\begin{corollary} \label{XII.1.2} For $F$ coherent given, and $m$ large enough, one has a canonical isomorphism
\steco
\begin{equation} \label{eq:XII.5} {}\H^i(X, F(-m))'\isomto\H^0(X, \SheafExt_{\OX}^{r-i}(F, \omx^r)(m))
\end{equation}
(where the $'$ denotes the dual vector space).
\end{corollary}
Indeed, on projective space $X$, one has for every pair of coherent sheaves $F, G$ and for $n$ large enough a canonical isomorphism:
\steco
\begin{equation} \label{eq:XII.6} {}\Ext^n(X;F(-m), G)\simeq \Ext^n(X;F, G(m)) \isomto \H^0(X, \SheafExt_{\OX}^{n}(F, G)(m)),
\end{equation}
(the isomorphism of the first two members being trivially true for every $m$), as follows from the spectral sequence of global $\Ext$
\[
\H^p(X, \SheafExt^q_{\OX}(F, G(m))) \To \Ext^{\boule}(X;F, G(m)),
\]
which degenerates for $m$ large enough thanks to the fact that
\[
\SheafExt^q_{\OX}(F, G(m)) \simeq \SheafExt^q_{\OX}(F, G)(m),
\]
and that the $\SheafExt^q_{\OX}(F, G)$ are coherent sheaves. Thus (\Ref{eq:XII.5}) follows from (\Ref{eq:XII.6}) and (\Ref{eq:XII.1}).

\begin{corollary} \label{XII.1.3}
For $i, F$ given, the following conditions are equivalent:
\begin{enumeratei}
\item
$\H^i(X, F(-m))=0$ for $m$ large.
\item[\textup{(i bis)}] $\H^i(X, F(\cdot))= \bigoplus_{m\in\ZZ}\H^i(X, F(-m))$ is an $S$-module of finite type, where $S=k[t_o, \ldots, t_r]$.
\item
$\SheafExt^{r-i}_{\OX}(F, \omx^r)=0$.
\item[\textup{(ii bis)}] $\SheafExt^{r-i}_{\OX}(F, \OX)=0$.
\item
$\H^i_x(F_x)=0$ for every closed point $x$ of $X$.
\item
$\H^{i+1}_x(\widetilde{F}_x)=0$ for every closed point $x$ of the punctured projecting cone $\widetilde{X}=\Spec(S)-\Spec(k)$ of $X$, where $\widetilde{F}$ denotes the inverse image of $F$ by the canonical morphism $\widetilde{X}\to X$.
\end{enumeratei}
\end{corollary}

\begin{proof}
(i) \SSI (i bis) since the submodule of $\H^i(X, F(\cdot))$ sum of the homogeneous components of degree $\geq \nu$ is of finite type over $S$ (in fact, for $i\neq 0$, it is even of finite type over $k$), (\cf FAC or \EGA III 2.2.1 and 2.3.2).

(i) \SSI (ii) by virtue of corollary \Ref{XII.1.2}.

(ii) \SSI (ii bis) since $\Omega_{X/k}^r$ is locally isomorphic to $\OX$.

(ii bis) \SSI (iii) by virtue of the local duality theorem for $\OXx$ (which is a regular local ring of dimension $r$), according to which the \og dual\fg of $\SheafExt^{r-i}_{\OX}(F, \OX)_x$ is identified with $\H^i_x(F_x)$ (V 2.1).

(ii bis) is equivalent to the analogous relation
\[
\SheafExt^{r-i}_{\Oo_{\widetilde{X}}}(\widetilde{F}, \Oo_{\widetilde{X}})=0
\]
(thanks to the fact that $\widetilde{X}\to X$ is faithfully flat, hence the inverse image of $\SheafExt^{r-i}_{\OX}(F, \OX)$ is isomorphic to $\SheafExt^{r-i}_{\Oo_{\widetilde{X}}}(\widetilde{F}, \Oo_{\widetilde{X}})$) and this last relation is equivalent to (iv) by the local duality theorem for the local ring $\OXx$, which is regular of dimension $r+1$.
\skipqed
\end{proof}

In particular, applying this to all $i\leq n$, one finds:

\begin{corollary} \label{XII.1.4}
Equivalent conditions for $n, F$ given:
\begin{enumeratei}
\item
$\H^i(X, F(-m))=0$ for $i\leq n$ and $m$ large.
\item[\textup{(i bis)}] $\H^i(X, F(\cdot))$ is an $S$-module of finite type for $i\leq n$.
\item
$\prof(F_x)>n$ for every closed point $x$ of $X$.
\item
$\prof(\widetilde{F}_x)>n+1$ for every closed point $x$ of $\widetilde{X}$.
\end{enumeratei}
\end{corollary}

The interest of corollaries \Ref{XII.1.3} and \Ref{XII.1.4} is to express a global condition (i) or~(i~bis) in terms of local conditions, namely the vanishing of local invariants such as $\H_x^i(X, F_x)$ or $\H_x^i(X, \widetilde{F}_x)$, or an inequality on the depth. In this form, these results remain trivially valid for an arbitrary projective scheme $X$, and a very ample invertible sheaf $\OX(1)$ on $X$, as one sees by inducing the latter with the help of a suitable projective immersion $X\to\PP_k^r$. (Of course, the conditions \Ref{XII.1.3}~(ii) and \Ref{XII.1.3}~(ii bis) are no longer equivalent to the others in this general case, except if one assumes that $X$ is regular for example). One may moreover generalize to the case of a projective morphism $X\to S$ in the following way:

\begin{proposition} \label{XII.1.5}
Let $f:X\to S$ be a projective morphism, with $S$ noetherian, $\OX(1)$ an invertible Module on $X$ very ample relative to $S$, $F$ a coherent Module on $X$, flat with respect to $S$, $s$ an element of $S$, $X_s$ the fibre of $X$ at $s$ (regarded as a projective scheme over $k(s)$), $F_s$ the sheaf induced on $X_s$ by $F$, finally $i$ an integer. Suppose that for every closed point $x$ of $X_s$, one has $\H^i_x(F_{s, x})=0$ (for example $\prof(F_{s, x})>i$). Then there exists an open neighbourhood $U$ of $s$, such that the same condition is satisfied for the $s'\in U$. Moreover, for such a $U$, one has
\[
\R^if_\ast(F(-m))= 0 \; \text{for}\; m \; \text{large, }
\]
and if $\Ss$ is a quasi-coherent graded algebra over $S$, generated by $\Ss^1$, and which defines~$X$ with $\OX(1)$ as $X=\Proj(\Ss)$, $\OX(1)=\Proj(\Ss(1))$, then the $S$-module
\[
\R^if_\ast(F(\cdot)) = \tbigoplus_{m\in\ZZ}\R^if_\ast(F(m))
\]
is of finite type over $U$.
\end{proposition}

Let us embed $X$ in an $X'=\PP_S^r$ in such a way that $\OX(1)$ is induced by $\Oo_{{X'}}(1)$ (this is possible provided one replaces $S$ by an affine neighbourhood of $s$). Let us set\refstepcounter{toto}\label{P5}\sfootnote{The first part of \Ref{XII.1.5} can be obtained at once by applying the purely local statement \EGA IV 12.3.4 to the preceding $\underline{E}^j$, which short-circuits the greater part of the proof that follows.} for every integer $j$, and every $t\in S$:
\steco
\begin{equation} \label{eq:XII.7} {}\underline{E}^j(t)= \SheafExt^j_{\Oo_{X'_t}} (F_{t'}, \Oo_{X'_t}(-r-1))
\end{equation}
Thus $\underline{E}^j(t)$ is a coherent Module on $X_t$, I say that for variable $t$, the family of these Modules is \og constructible\fg in the following sense: there exists for every $t\in S$ a non-empty open $V$ of the closure $\overline{t}$, which one endows with the induced reduced structure, and a coherent Module $\underline{E}^j(V)$ on $X_V=X\times_SV$, \emph{flat} relative to $V$, such that for every $t'\in V$, $\underline{E}^j(t)$ is isomorphic to the Module induced by $\underline{E}^j(V)$ on $X_t$. To verify this assertion, setting $Z=\overline{t}$ endowed with the induced structure, one considers the coherent Modules
\[
\underline{E}^j(Z)= \SheafExt^j_{\Oo_{X'_Z}} (F_Z, \Oo_{X'_Z}(-r-1))
\]
(where the index $Z$ still means that one induces over $Z$), and one takes for $V$ a non-empty open of $Z$ such that the Modules $\underline{E}^j(Z)$ are flat over $V$: this is possible, since one verifies at once that $\underline{E}^j(Z)=0$ for $j$ not belonging to the interval $[0, r]$, and one may apply \SGA 1 IV 6.11. One then takes for $\underline{E}^j(V)=\underline{E}^j(Z)|_{X_V}$, and one verifies easily that it answers the question.

From the preceding remark it follows that there exists a finite partition of $S$ formed by sets $V_\alpha$ of the form $V=V(t)$ as above, (noetherian induction), and applying Serre's theorem \EGA III 2.2.1 to the $\underline{E}^j(V_i)$, one sees that there exists an integer $m_0$ such that
\[
\R^if_{V_\alpha\ast}(\underline{E}^j(V_\alpha))=0\textup{ for } i\neq 0, m\geq m_0, \textup{ for }j,
\]
from which it follows, using the flatness of $\underline{E}^j(V_\alpha)$ with respect to $V_\alpha$ and easy relations \`a la Künneth (\cf \EGA III par\ptbl7), that
\[
\H^i(X_t, \underline{E}^j(t)(m))=0\textup{ for } i\neq 0, m\geq m_0, \textup{ for }j,
\]
for every $t\in V_\alpha$, hence for every $t$ since the $V_\alpha$ cover $S$. From this and from the spectral sequence of global Ext there follows, thanks to \Ref{XII.1.1}, as in the proof of \Ref{XII.1.2}, an isomorphism
\steco
\begin{equation} \label{eq:XII.8} {}\H^i(X_t, F_t(-m))' \isomfrom \H^0(X_t, \underline{E}^{r-i}(t)(m))\textup{ for } m\geq m_0,
\end{equation}
every integer $i$, and every $t\in S$.

Let us now use the hypothesis made on $F_s$, which is written
\steco
\begin{equation} \label{eq:XII.9} {}\underline{E}^{r-i}(s)=0,
\end{equation}
and thanks to (\Ref{eq:XII.8}) is equivalent to
\steco
\begin{equation} \label{eq:XII.10} {}\H^i(X_s, F_s(-m))=0 \; \text{for}\; m\geq m_0.
\end{equation}
As $F$ hence $F(-m)$ is flat with respect to $S$, it follows by
the relations \`a la Künneth already invoked, that (for given $m$) the
same relation (\Ref{eq:XII.10}) holds upon replacing $s$ by a
$t$ neighbouring $s$, in particular for every generalization $t$ of
$s$. By virtue of (\Ref{eq:XII.8}), one will therefore have, for such a
generalization \steco
\begin{equation} \label{eq:XII.11} {}\underline{E}^{r-i}(t)=0,
\end{equation}
now the set of $t\in s$ for which this relation is true is evidently a constructible set (for it induces on each $V_\alpha$ an open set); as it contains the generalizations of $s$, it contains an open neighbourhood $U$ of $s$. This proves the first assertion of \Ref{XII.1.5}. Moreover, for $t\in U$, one concludes from (\Ref{eq:XII.11}) and (\Ref{eq:XII.8}) that
\steco
\begin{equation} \label{eq:XII.12} {}\H^i(X_t, F_t(-m)=0 \; \text{for}\; m\geq m_0, \; t\in U,
\end{equation}
which by virtue of the relations \`a la Künneth, implies (in fact, is much stronger than)
\steco
\begin{equation} \label{eq:XII.13} {}\R^if_\ast(F(-m))=0 \; \text{on}\; U, \; \text{for} \; m\geq m_0.
\end{equation}
This proves the second assertion of \Ref{XII.1.5}. Finally the last follows from it at once, by proceeding as at the beginning of the proof of \Ref{XII.1.3}.

\refstepcounter{subsection}\label{XII.1.6}
\begin{enonce*}[remark]{Remark 1.6\sfootnotemark} \sfootnotetext{This remark is made more precise by the footnote on page~\pageref{P5}.} The proof simplifies considerably (eliminating any consideration of constructibility) when one already assumes that the hypothesis made for $s\in S$ is satisfied at all the $s'\in S$. In fact, when one makes the hypothesis that $F_s$ is of depth $>i$ at the closed points of $X$, one has a general statement, \emph{of a local nature on} $X$, which says that the same condition is satisfied for all the $X_t$, provided one replaces $X$ by a suitable open neighbourhood of the fibre $X_s$ (in other words, a certain subset of $X$, defined by conditions on the Modules induced by $F$ on the fibres, is open, \cf \EGA IV). As $f$ is here proper, one will therefore be able to take this neighbourhood of the form $f^{-1}(U)$, which recovers the first assertion of \Ref{XII.1.5} without painful d\'evissage. In this general case, one can still prove by the method of \loccit that the first assertion of \Ref{XII.1.5} (proved here by a global method, using that $X$ is projective over $S$) still follows from a purely local statement on $X$ (which the reader will make explicit if he finds it useful).
\end{enonce*}

% English body fragment for SGA 2 XII (en-2.tex)
\section[Lefschetz theory for a projective morphism: comparison]
{Lefschetz theory for a projective morphism: Grauert's comparison theorem} \label{XII.2}

\setcounter{equation}{13}
It is the following theorem:

\begin{theorem} \label{XII.2.1}
Let $f:X\to S$ be a projective morphism, with $S$ noetherian, $\OX(1)$ an invertible Module on $X$, ample relative to $S$, $Y$ the prescheme of zeros of a section $t$ of $\OX(1)$, $J$ the ideal defining $Y$, $X_n$ the subprescheme of $X$ defined by $J^{n+1}$, $\widehat{X}$ the formal completion of $X$ along $Y$, $\widehat{f}:\widehat{X}\to S$ the composite morphism $\widehat{X}\to X\to S$, $F$ a coherent Module on $X$, flat relative to $S$. One assumes moreover that for every $s\in S$, the Module $F_s$ induced on the fibre $X_s$ is of depth $>n$ at the points of the said fibre, and that $t$ is $F$-regular. Under these conditions:
\begin{enumeratei}
\item
The canonical homomorphism
\[
\R^if_\ast(F) \to \R^i\widehat{f}_\ast(\widehat{F})
\]
is an isomorphism for $i<n$, a monomorphism for $i=n$.
\item
The canonical homomorphism
\[
\R^i\widehat{f}_\ast(\widehat{F}) \to \varprojlim_m \R^if_\ast(F_m)
\]
is an isomorphism for $i\leq n$.
\end{enumeratei}
\end{theorem}

\begin{proof}
One immediately reduces to the case where $S$ is affine, and then to proving the

\begin{corollary} \label{XII.2.2}
Under the conditions of \Ref{XII.2.1} suppose moreover that $S$ is affine. then:
\begin{enumeratei}
\item
The canonical homomorphism
\[
\H^i(X, F) \to \H^i(\widehat{X}, \widehat{F})
\]
is an isomorphism for $i<n$, a monomorphism for $i=n$.
\item
The canonical homomorphism
\[
\H^i(\widehat{X}, \widehat{F}) \to \varprojlim_m \H^i(X_m, F_m)
\]
is an isomorphism for $i\leq n$.
\end{enumeratei}
\end{corollary}

Replacing if need be $\OX(1)$ by a tensor power, and $t$ by a power of~$t$, one may assume $\OX(1)$ very ample relative to $S$. On the other hand $t$ and hence $t^m$ being $F$-regular, multiplication by $t^m$, regarded as a homomorphism from $F(-m)$ into $F$, is injective, hence one has for every $m\geq 0$ an exact sequence:
\steco
\begin{equation} \label{eq:XII.14}
0\to F(-m) \lto{t^m} F \to F_m \to0,
\end{equation}
whence an exact sequence of cohomology
\[
\H^i(X, F(-m)) \to \H^i(X, F) \to \H^i(X, F_m) \to \H^{i+1}(X, F(-m)).
\]
Now by virtue of \Ref{XII.1.5} one has $\H^i(X, F(-m))=0$ for $i\leq n$, and $m$ sufficiently large, which proves the

\begin{lemma} \label{XII.2.3}
For $m$ large, the canonical homomorphism
\[
\H^i(X, F) \to \H^i(X, F_m)
\]
is bijective if $i<n$, injective if $i=n$.
\end{lemma}

%
This shows that for $i<n$, the projective system $(\H^i(X_m, F_m))_{m\geq 0}$ is essentially constant, \afortiori satisfies the Mittag-Leffler condition, hence (taking account of $\widehat{F}= \varprojlim F_m$) one concludes (ii) by $\EGA 0_{\textup{III}}$ 13.3. On the other hand (i) follows trivially from this, taking account of \Ref{XII.2.3}.

\refstepcounter{subsection}\label{XII.2.4}
\begin{enonce*}{Corollary 2.4\ndemark}\ndetext{see Corollary I.1.4 of the article of Mme Raynaud (Raynaud~M., {\og Th\'eor\`emes de Lefschetz en cohomologie des faisceaux coh\'erents et en cohomologie \'etale. Application au groupe fondamental\fg}, \emph{Ann. Sci. \'Ec. Norm. Sup. (4)} \textbf{7} (1974), p\ptbl 29--52).}
Let $f:X\to S$ be a projective and flat morphism, with $S$ locally noetherian, $\OX(1)$ an invertible Module on $X$, ample relative to $S$, $t$ a section of this Module which is $\OX$-regular, $Y$ the subprescheme of zeros of $t$, $\widehat{X}$ the formal completion of $X$ along $Y$. One assumes that for every $s\in S$, $X_s$ is of depth $\geq 1$ (\resp of depth $\geq 2$) at its closed points. Then for every open neighbourhood $U$ of $Y$, the functor
\[
F \mto \widehat{F}
\]
from the category of coherent locally free Modules on $U$, to the category of coherent locally free Modules on $\widehat{X}$, is faithful (\resp fully faithful, \ie the Lefschetz condition ($\Lef$) of \Ref{X}~\Ref{X.2} is satisfied).
\end{enonce*}
Introducing for two locally free Modules $F$ and $G$ on $U$ the Module
\[
H=\SheafHom_{\Oo_U}(F, G)
\]
one is reduced to proving that the canonical homomorphism
\steco
\begin{equation} \label{eq:XII.15} {}\H^0(U, H) \to \H^0(\widehat{U}, \widehat{H})
\end{equation}
is injective (\resp bijective). Now the Modules $H_t$ are of depth $\geq 1$ (\resp $\geq 2$) at the closed points of $X_t$, one can therefore apply \Ref{XII.2.1}, which implies the conclusion \Ref{XII.2.4} in the case where $U=X$. In the case of an arbitrary $U$, one notes that the question is local on $S$, hence one may assume $S$ affine. Then every coherent Module on $X$ is a quotient of a coherent locally free Module, (since $\OX(1)$ is a relatively ample invertible Module on $X$). As the dual Module $H'=\SheafHom(H, \Oo_U)$ extends to a coherent Module on $X$, which is therefore isomorphic to a cokernel of a homomorphism of locally free Modules on $X$, it follows by transposition that one can find a homomorphism
\[
u':{L'}^0 \to {L'}^1
\]
of locally free Modules on $X$, inducing a homomorphism
\[
u:L^0 \to L^1
\]
of locally free Modules on $U$, such that one has an exact sequence
\[
0\to H \to L^0 \lto{u} L^1.
\]
Using the five lemma (which becomes the three lemma), and the left exactness of the functor $H^0$, one is reduced to proving that (\Ref{eq:XII.15}) is injective (\resp bijective) when one replaces $H$ therein by $L^0$, $L^1$, which reduces us to the case where $H$ is induced by a locally free Module $H'$ on $X$. Moreover in the non-resp.\ case this reduction is even unnecessary, for the kernel of (\Ref{eq:XII.15}) is in any case formed of the sections of $H$ on~$U$ which vanish in a suitable open neighbourhood $V$ of $Y$, now the restriction homomorphism $\H^0(U, H)\to\H^0(V, H)$ is injective, for $H$ is of depth $\geq 1$ at the points of every closed subset $Z$ of $X$ not meeting $Y$ (\cf lemma below). In the resp.\ case, one is reduced to proving that
\[
H^0(X, H')\to\H^0(U, H')
\]
is bijective, which follows from the fact that $H'$ is of depth $\geq 2$ at all the points of a closed set $Z=X- U$ of $X$ not meeting $Y$. One must therefore simply prove the

\begin{lemma} \label{XII.2.5} Let $F$ be a coherent Module on $X$, flat with respect to $S$, such that for every $s\in S$, $F_s$ is of depth $\geq n$ at every closed point of $X_s$. Then for every closed subset $Z$ of $X$ not meeting $Y$, $F$ is of depth $\geq n$ at all the points of $Z$.
\end{lemma}

Indeed, for every $x\in X$, setting $s=f(x)$, one has
\steco
\begin{equation} \label{eq:XII.16} {}\prof(F_x)\geq \prof(F_{s, x}),
\end{equation}
as one sees by lifting in any manner a maximal $F_{s, x}$-regular sequence of elements of $\rr(\Oo_{X_s, x})$, which gives an $F_x$-regular sequence by virtue of \SGA 1 IV 5.7. Now if $x$ belongs to a $Z$ as in lemma \Ref{XII.2.5}, then $x$ is necessarily closed in $X_s$, in other words, $Z$ is \emph{finite} over $S$. Indeed $Z$ (equipped with a structure induced by $X$) is projective over $S$ as a closed subprescheme of $X$ which is so, and $Z$ is affine over $S$ as a closed subprescheme of $X- Y$, which is so.
\skipqed
\end{proof}

\begin{remark} \label{XII.2.6}
Suppose that for every $s\in S$, the section $t_s$ of $\Oo_{X_s}(1)$ induced by $t$ is $\Oo_{X_s}$-regular (which implies by \SGA 1 IV 5.7 that $t$ is $\OX$-regular). Then the hypotheses made are stable under extension of the base $S'\to S$ ($S'$ locally noetherian). Hence the conclusion remains valid after any change of base.
\end{remark}

\section[Lefschetz theory for a projective morphism: existence]{Lefschetz theory for a projective morphism: existence theorem}\label{XII.3}%

\setcounter{equation}{16}%
\refstepcounter{subsection}\label{XII.3.1}%
\begin{enonce*}{Theorem 3.1\ndemark}
\ndetext{for a version without a flatness hypothesis, see (Raynaud~M., {\og Th\'eor\`emes de Lefschetz en cohomologie des faisceaux coh\'erents et en cohomologie \'etale. Application au groupe fondamental\fg}, \emph{Ann. Sci. \'Ec. Norm. Sup. (4)} \textbf{7} (1974), p\ptbl 29--52, theorem II.3.3).}
Let $f:X\to S$ be a projective morphism, with $S$ noetherian, $\OX(1)$ an invertible Module on $X$ ample relative to $S$, $X_0$ the subprescheme of zeros of a section $t$ of $\OX(1)$, $\widehat{X}$ the formal completion of $X$ along $X_0$, $\formelF$ a coherent Module on $\widehat{X}$, $F_0$ the Module it induces on $X_0$. One assumes moreover:
\begin{enumeratea}
\item\label{XII.3.1a}
$\formelF$ is flat with respect to $S$.
\item\label{XII.3.1b}
For every $s\in S$, the section $t_s$ induced by $t$ on $X_s$ is $\formelF_s$-regular (which implies that $F_0$ is likewise flat with respect to $S$, \cf \textup{\SGA 1 IV 5.7}).
\item\label{XII.3.1c}
For every $s\in S$, $F_{0s}$ is of depth $\geq 2$ at the closed points of $X_{0s}$.
\end{enumeratea}
One assumes moreover that $S$ admits an ample invertible sheaf. Under these conditions, there exists a coherent Module $F$ on $X$, and an isomorphism of its formal completion $\widehat{F}$ with $\formelF$.
\end{enonce*}

This statement will follow from the following:

\begin{corollary} \label{XII.3.2}
Under the conditions \Ref{XII.3.1a}), \Ref{XII.3.1b}), \Ref{XII.3.1c}) above one has the following:
\begin{enumeratei}
\item
The Module $\widehat{f}_\ast(\formelF)$ on $S$ is coherent, hence for every $n$, the Module $\widehat{f}_\ast(\formelF(n))$ on~$S$ is coherent.
\item
For $n$ large, the canonical homomorphism $\widehat{f}^\ast\widehat{f}_\ast(\formelF(n))\to \formelF(n)$ is surjective.
\end{enumeratei}
\end{corollary}

Let us admit the corollary for the moment, and prove \Ref{XII.3.1}. Thanks to the last hypothesis made in \Ref{XII.3.1}, one can reduce to the case where $X=\PP^r_S$, replacing if need be $\OX(1)$, $t$ by a suitable power. I say that one may assume moreover that for every~$s$, one has $t_s\neq 0$. Otherwise, one has indeed $\formelF_s=0$ by b), or what comes to the same by Nakayama, $F_{0s}=0$ \ie $s$ does not belong to the image of $\Supp F_0$ under the morphism $f_0:X_0\to S$ induced by $f$. Now this image $S'$ is open by virtue of a), b) since $F_0$ is flat with respect to $S$, and it is evident that it suffices to prove the conclusion of \Ref{XII.3.1} in the situation obtained by restricting above $S'$, for the coherent Module $F'$ on $X|_{S'}$ obtained by will be the restriction of a coherent Module $F$ on $X$, which will answer the question. One may therefore assume that in addition to the hypotheses a), b), c), the following hypotheses are likewise satisfied:
\begin{enumerate}
\item[\textup{a$'$)}] $\OX$ is flat with respect to $S$.
\item[\textup{b$'$)}] For every $s\in S$, the section $t_s$ is $\Oo_{X_s}$-regular.
\item[\textup{c$'$)}] For every $s\in S$, $\Oo_{X_{0s}}$ is of depth $\geq 2$ at the closed points of $X_{0s}$. (It~suffices to choose $X=\PP_S^r$ with $r\geq 3$, which is permissible).
\end{enumerate}
Now \Ref{XII.3.2} implies that one can find an epimorphism
\steco
\begin{equation} \label{eq:XII.17}
\widehat{L} \to \formelF \to 0,
\end{equation}
where $L$ is a Module on $X$ of the form $f_\ast(G)(-n)$, $G$ being a coherent locally free Module on $S$: indeed it suffices; for $n$ large, to represent the coherent Module $\widehat{f}_\ast(\formelF)$ on $S$ as a quotient of such a $G$. On the other hand, the hypotheses a), b), c) on $f$, $t$ imply that $\widehat{L}$ satisfies the same conditions a), b), c) as $\formelF$. One easily concludes that the same holds for the kernel of (\Ref{eq:XII.17}), to which one can therefore apply the same argument, so that $\formelF$ is represented as the cokernel of a homomorphism
\steco
\begin{equation} \label{eq:XII.18}
\widehat{L'} \to \widehat{L},
\end{equation}
where $L$, $L'$ are locally free Modules on $X$. Now by virtue of a$'$) and the second part of c$'$), and of \Ref{XII.2.1} or \Ref{XII.2.4} at choice, the homomorphism (\Ref{eq:XII.18}) comes from a homomorphism $L'\to L$ of Modules on $X$. It now suffices to take for $F$ the cokernel of $L'\to L$, and one is done.

It remains to prove \Ref{XII.3.2}. This had been done in the
seminar by a somewhat painful expedient, consisting in
interpreting everything in terms of cohomology on the punctured
projecting cone of $X$ relative to $S$, in order to be able to
reduce to the theorem \Ref{IX.2.1}. A more direct and more
satisfactory way (although substantially identical), seems to
me now the following. It consists in noting that in \Ref{IX},
\numero \Ref{IX.2}\refstepcounter{toto}\label{XII.14} (and with the
notations of that expos\'e) the hypothesis that the morphism
$f:\Xx\to\Xx'$ be adic does not intervene anywhere in the
proof of \Ref{IX.2.1}, via $\EGA 0_{\textup{III}}$ 13.7.7;
it suffices to assume instead that $\Xx$ is likewise adic,
and to choose two ideals of definition $\Jj$ for $\Xx'$, $\Ii$
for $\Xx$, such that $\Jj\Oo_\Xx\subset\Ii$, and to define
$\Ss=\gr_\Jj(\Oo_{\Xx'})$, and to consider $\gr_\Ii(\formelF)$. In
any case, \Ref{IX.2.1} can be applied directly to the
morphism $\widehat{f}:\widehat{X}\to S$ considered in the present
no., where one simply takes $\Jj=0$. Thus to verify that
$\widehat{f}_\ast(\formelF)$ is coherent, it suffices by virtue of
\loccit to verify that $\R^if_{0\ast}(\gr_\Ii(\formelF))$ is
coherent on $S$ for $i=0, 1$; for this one notes that by virtue of
a) and b), the Module under consideration is none other than
$\bigoplus_{m\geq 0}\R^if_{0\ast}(F_0(-m))$, which
is indeed coherent by virtue of hypothesis c) and of
\Ref{XII.1.5}.

This proves \Ref{XII.3.2} (i). For \Ref{XII.3.2} (ii), we will need the

\begin{lemma} \label{XII.3.3}
Under the conditions of \Ref{XII.3.1a}), \Ref{XII.3.1b}), \Ref{XII.3.1c}) of \Ref{XII.3.1}, set
\[
G_m= \widehat{f}_\ast(\formelF(\cdot)_m) = \tbigoplus_n\widehat{f}_\ast(\formelF_m(n))
\]
Then the projective system $(G_m)$ satisfies the Mittag-Leffler condition.
\end{lemma}

One may assume $S$ affine, with ring $A$. Let then $\Ss$ be a graded $A$-algebra of finite type with positive degrees, and $t'\in\Ss_1$, such that $X$ immerses in $\Proj(\Ss)$, $\OX(1)$ being induced by $\Oo(1)$ and the section $t$ being the image of $t'$. Let us equip $\Ss$ with the $\Jj$-adic filtration, where $\Jj=t'\Ss$, and consider the projective system of the $\formelF(\cdot)_m$ in the category of abelian sheaves on $X_0$. One is still under the preliminary conditions of $\EGA 0_{\textup{III}}$ 13.7.7\sfootnote{Rectified as indicated in \Ref{IX} p\ptbl\pageref{pagerectif}.} and moreover $\H^i(X_0, \gr(\formelF(\cdot))$ is a Module of finite type over $\gr_\Jj(\Ss)$ for $i=0, 1$. Indeed, as $t'$ is $\formelF$-regular, one sees immediately that as a Module over $(\Ss /t'\Ss)[T]$ (of which $\gr_\Jj(\Ss)$ is a quotient), the Module under consideration identifies with $\H^i(X_0, F_0(\cdot))\otimes_{\Ss /t'\Ss}(\Ss /t'\Ss)[T]$, now by virtue of \Ref{XII.1.5}, $\H^i(X_0, F_0(\cdot))$ is of finite type over $\Ss$, hence over $\Ss /t'\Ss$, for $i=0, 1$, which proves our assertion. Consequently one is under the conditions of application of $0_{\text{III}}$ 13.7.7 with $n=1$, which proves \Ref{XII.3.3}.

This point being granted (and still assuming $S$ affine, which is permissible in order to prove \Ref{XII.3.2} (ii)) let $m_0$ be such that $m\geq m_0$ implies $\Im(G_m\to G_0)= \Im(G_{m_0}\to G_0)$, so that both sides are also equal to $\Im(\varprojlim G_m\to G_0) = \Im(\widehat{f}_\ast(\formelF(\cdot))\to f_\ast\formelF_0(\cdot))$. Let us now remark that for $n$ large, $\formelF_{m_0}(n)$ is generated by its sections, hence $\formelF_0(n)$ is generated by sections which lift to $\formelF_{m_0}$, hence (thanks to the choice of~$m_0$) which lift to $\formelF$. Hence the sections of $\formelF$ generate $\formelF_0$, hence also~$\formelF$ thanks to Nakayama. This proves \Ref{XII.3.2} (ii), hence \Ref{XII.3.1}.

\begin{corollary} \label{XII.3.4}
Let $f:X\to S$ be a projective and flat morphism, with $S$ locally noetherian, $\OX(1)$ an invertible Module on $X$, ample relative to $S$, $t$ a section of this Module such that for every $s\in S$, the section $t_s$ induced on the fibre $X_s$ is $\Oo_{X_s}$-regular, $X_0$ the subprescheme of zeros of $t$, $\widehat{X}$ the formal completion of $X$ along~$X_0$. One assumes that for every $s\in S$, $X_{0s}$ is of depth $\geq 2$ at its closed points, (\ie $X_s$ is of depth $\geq 3$ at the closed points of $X_{0s}$) and $X_s$ is of depth $\geq 2$ at its closed points. Under these conditions, the pair $(X, X_0)$ satisfies the effective Lefschetz condition ($\Leff$) of paragraph~\Ref{X.2} of expos\'e~\emph{\Ref{X}}, \ie:
\begin{enumeratea}
\item
For every open neighbourhood $U$ of $X_0$ in $X$, the functor
\[
F \mto \widehat{F}
\]
from the category of coherent locally free Modules on $U$ to the category of coherent locally free Modules on $\widehat{X}$ is fully faithful.
\item
For every coherent locally free Module $\formelF$ on $\widehat{X}$, there exists an open neighbourhood~$U$ of $X_0$, and a coherent locally free Module $F$ on $U$ such that $\formelF$ is isomorphic to~$\widehat{F}$.
\end{enumeratea}
\end{corollary}

Indeed, a) has already been noted in \Ref{XII.2.4} under weaker conditions. For b), one applies \Ref{XII.3.1} which gives the conclusion, at least if $S$ is noetherian and admits an absolutely ample invertible Module, in particular if $S$ is affine. Indeed, if $F$ is a coherent module on $X$ such that $\widehat{F}$ is isomorphic to $\formelF$ hence locally free, it follows that $F$ is locally free on a neighbourhood $U$ of $X_0$, and $F|_U$ will satisfy the desired condition. But let us now note that by virtue of \Ref{XII.2.5}, for such an $F$, its image under the immersion $U\to X$ is coherent, and moreover independent of the chosen solution $(U, F)$ (taking account of the fact that two solutions coincide in a neighbourhood of $X_0$, by virtue of a)). More precisely, one can find a coherent Module $F$ on $X$ and an isomorphism $\widehat{F}\isomto \formelF$, such that $F$ is of depth $\geq 2$ at all the points of $X$ which are closed in their fibre, and this determines $F$ up to a unique isomorphism. Thanks to this uniqueness property, the solutions of the problem that one finds by inducing above the affine open sets of $S$ glue together, whence a coherent $F$ on the whole of $X$ and an isomorphism $\widehat{F}\simeq \formelF$. Restricting $F$ to the open set $U$ of the points at which it is free, one finds what one sought.

Thanks to \Ref{XII.2.4} and \Ref{XII.3.4}, one can exploit, in the situation of a projective algebraic scheme and a ``hyperplane section'' of the same, the general facts established in expos\'es \Ref{X} and \Ref{XI} concerning the conditions (Lef) and (Leff). Thus:

\begin{corollary} \label{XII.3.5} Let $X$ be a projective algebraic scheme equipped with an ample invertible Module $\OX(1)$, let $t$ be a section of this Module which is $\OX$-regular, and let $X_0$ be the subscheme of zeros of $t$. Suppose that $X$ is of depth $\geq 2$ at its closed points (\resp and of depth $\geq 3$ at the closed points of $X_0$). Then $\pi_0(X_0)\to \pi_0(X)$ is bijective, in particular $X$ is connected if and only if $X_0$ is, and choosing a geometric base-point in $X_0$, $\pi_1(X_0)\to \pi_1(X)$ is surjective, and more generally for every open set $U\supset X_0$, the homomorphism $\pi_1(X_0)\to \pi_1(U)$ is surjective (\resp the homomorphism $\pi_1(X_0)\to \varprojlim_U\pi_1(U)$ is bijective). In the resp.\ case, if one assumes moreover that the local ring of every closed point of $X$ not in $X_0$ is pure (\Ref{X.3.2}) (for example is regular, or merely a complete intersection) then $\pi_1(X_0)\to \pi_1(X)$ is an isomorphism.
\end{corollary}

One applies \Ref{X.2.4} and \Ref{XII.3.4}. One will note that in the resp.\ case, the hypothesis that $X$ be of depth $\geq 3$ at the closed points of $X_0$, implies that all the irreducible components of dimension $\neq 0$ of $X$ are of dimension $\geq 3$ (as one sees by noting that such a component necessarily meets $X_0$, and looking at a closed point of the intersection).

\begin{remarkstar}
When $X$ is normal, of dimension $\geq 2$ at all its points, it is of depth $\geq 2$ at its closed points and one is under the non-resp.\ conditions of \Ref{XII.3.5}. In this case, one has a more elementary proof of the surjectivity of $\pi_1(X_0)\to \pi_1(X)$ with the aid of Bertini's theorem (\cf \SGA 1 X.2.10). When one assumes moreover $X_0$ normal, and $X$ of dimension $\geq 3$ at all its points, then one is under the resp.\ conditions of \Ref{XII.3.5}. In this case, \Ref{XII.3.5} was established by Grauert (it turns out indeed that thanks to the normality hypothesis, one then manages to do without the existence theorem \Ref{XII.3.1} by expedients). It is this proof of Grauert which was the starting point of the ``Lefschetz theory'' which is the object of the present seminar.
\end{remarkstar}

\begin{corollary} \label{XII.3.6} Let $X$, $\OX(1)$, $t$, $X_0$ be as in \Ref{XII.3.5}. Suppose that $X$ is of depth $\geq 2$ at its closed points, and that
\[
\H^i(X_0, \Oo_{X_0}(-n))=0
\]
for $n>0$ and for $i=1$ (\resp for $i=1$ and for $i=2$), which implies by virtue of \Ref{XII.1.4} that $X_0$ is of depth $\geq 2$ (\resp $\geq 3$) at its closed points, \ie that $X$ is of depth $\geq 3$ (\resp $\geq 4$) at the closed points of $X_0$. Under these conditions, for every open neighbourhood $U$ of $X_0$, $\Pic(U)\to\Pic(X_0)$ is injective, in particular $\Pic(X)\to\Pic(X_0)$ is injective, (\resp $\varprojlim_U \Pic(U)\to\Pic(X_0)$ is bijective). In the resp.\ case, if one assumes moreover that the local ring of $X$ at every closed point not in $X_0$ is parafactorial (\Ref{XI.3.1}) (for example is regular, or more generally a complete intersection), then $\Pic(X)\to\Pic(X_0)$ is bijective.
\end{corollary}
One applies \Ref{XI}~\Ref{XI.3.12} and \Ref{XI.3.13}, noting that the resp.\ hypothesis implies that the irreducible components of dimension $\neq 0$ of $X$ are of dimension $\geq 4$. One finds in particular, applying this to the case where $X$ is a global complete intersection of dimension $\geq 4$ in projective space:

\begin{corollary} \label{XII.3.7} Let $X$ be an algebraic scheme of dimension $\geq 3$, which is a complete intersection in a scheme $\PP_k^r$. Then $\Pic(X)$ is the free group generated by the class of the sheaf $\OX(1)$.
\end{corollary}
One reasons by induction on the number of hypersurfaces of which $X$ is the intersection, applying \Ref{XII.3.6} and noting that for a complete intersection $X$ of dimension $\geq 3$, one $\H^i(X, \OX(n))=0$ for $i=1, 2$, and every $n$.

\begin{remark} \label{XII.3.8}
In the case where $X$ is a nonsingular hypersurface, \Ref{XII.3.7} is due to Andreotti. The result \Ref{XII.3.7} is also expressed (when $X$ is nonsingular) by saying that the homogeneous coordinate ring of $X$ is factorial, and in this form is contained in \Ref{XI}~\Ref{XI.3.13} (ii). Let us also point out that Serre had given a proof of \Ref{XII.3.7} in the nonsingular case, by transcendental means, using a specialization argument to reduce to the case of characteristic~$0$, where one has the Lefschetz theorem in its classical form. Of course, the fact that the purely algebraic proof given here allows one to rid oneself of nonsingularity hypotheses in the statement of the Lefschetz theorem, invites one to reconsider the latter equally in the classical case. \Cf the following expos\'e which proposes conjectures in this sense.

In the corollaries \Ref{XII.3.5} and \Ref{XII.3.7} we have placed ourselves over a base field, whereas the key theorems \Ref{XII.2.4} and \Ref{XII.3.4} are valid over an arbitrary base. To generalize to a general $S$ these corollaries \Ref{XII.3.5} and \Ref{XII.3.6}, we must give serviceable criteria in order that a point of $X$ (flat over $S$) have a ``pure'' local ring \resp a parafactorial one. This will be the object of the following no.
\end{remark}

\section{Formal completion and normal flatness} \label{XII.4}

\setcounter{equation}{18}

\begin{theorem} \label{XII.4.1} Let $X$ be a locally noetherian prescheme, locally immersible in a regular scheme, $Y$ a closed subset of $X$, $U=X- Y$, $X_0$ a closed subprescheme of $X$ defined by an ideal $\Jj$, $\widehat{X}$ the formal completion of $X$ along $X_0$, $U_0$ the trace of $X_0$ on $U$, $\widehat{U}$ the formal completion of $U$ along $U_0$, $i\colon U\to X$ and $\widehat{i}\colon \widehat{U}\to \widehat{X}$ the canonical immersions, $n$ an integer. Assume
\begin{enumeratea}
\item
$X$ is normally flat along $X_0$ at the points of $Y\cap X_0$, \ie at these points the Modules $\Jj^n/\Jj^{n+1}$ on $X_0$ are flat \ie locally free.
\item
For every $x\in Y\cap X_0$, one has $\prof \Oo_{X_0, x}\geq n+2$.
\end{enumeratea}
Under these conditions, one has the following:
\begin{enumerate}
\item Let $F$ be a coherent Module on $U$, suppose that one has:
\begin{enumerate}
\item[\textup{c)}] For every $x\in Y- Y\cap X_0$, one has $\prof \Oo_{X, x}\geq n+2$.

\item[\textup{d)}] $F$ is free at the points of $U_0$, and of depth $\geq n+1$ at all the points of $U$ where it is not free.
\end{enumerate}
Then the graded Module
\[
\tbigoplus_{m\geq 0}\R^pi_\ast(\Jj^mF)
\]
on $\bigoplus_{m\geq 0}\Jj^m$ is of finite type for $p\leq n$. \item Let $\formelF$ be a coherent Module on $\widehat{U}$, then the graded Module
\[
\tbigoplus_{m\geq 0}\R^pi_\ast(\Jj^m\formelF/\Jj^{m+1}\formelF)
\]
on $\bigoplus_{m\geq 0}\Jj^m/\Jj^{m+1}=\gr_{\Jj}(\OX)$ is of finite type for $p\leq n$.
\end{enumerate}
\end{theorem}

\begin{proof}
1) Let $X'\!=\!\Spec(\bigoplus_{m\geq 0}\Jj^m)$; the change of base \hbox{$f:X'\!\to\! X$} then defines $U'=X'- Y'$, $X_0'$, $U_0'=X_0' U'$, and immersions $i'\colon U'\to X'$, \hbox{$i_0':U_0'\to X_0'$}. One thus has a cartesian square
\[
\xymatrix{
X'\ar[d]_f&\ar[l]_-{i'} U'\ar[d]^g \\
X&\ar[l]_-{i} U
}
\]
and one has
\steco
\begin{equation} \label{eq:XII.19}
\tbigoplus_{m\geq 0}\R^pi_\ast(\Jj^mF) = \R^pi_\ast \Big(\tbigoplus_{m\geq 0} \Jj^mF \Big) =\R^pi_\ast (g_\ast(F')),
\end{equation}
where $F'=g^\ast F$, so that one does indeed have a canonical isomorphism
\[
g_\ast(F') \isomto \tbigoplus_{m\geq 0}\Jj^mF,
\]
for this is true at the points of $U_0$, from the fact that $F$ is free there by virtue of d), and equally at the points of $U_0$, from the fact that one has $\Jj^m=\Oo_U$ there, (so that in both cases, $\Jj^m\otimes_{\Oo_U}F\to \Jj^mF$ is an isomorphism).

On the other hand, as $f$ and consequently $g$ are affine, one has
\steco
\begin{equation} \label{eq:XII.20} {}\R^pi_\ast(g_\ast(F')) = \R^p(ig)_\ast(F') = \R^p(fi')_\ast(F') = f_\ast(R^pi'_\ast(F'))
\end{equation}
hence comparing (\Ref{eq:XII.19}) and (\Ref{eq:XII.20}), one sees that assertion 1) is equivalent to the following: $R^pi'_\ast(F')$ is a Module of finite type \ie coherent on $X'$, for every $p\leq n$. Now as $X$ is locally immersible in a regular scheme, the same is true of $X'$ which is of finite type over $X$, and one may apply the coherence criterion \Ref{VIII}~\Ref{VIII.2.3} to a coherent extension $F^{''}$ of $F'$: one wishes to express that $\SheafH^p_Y(F^{''})$ is coherent for $p\leq n+1$, and this is also equivalent to saying that for every $x'\in U'$ such that
\begin{equation*} \label{eq:XII.20bis} \tag{20 bis} {}\codim (\overline{x'}\cap Y', \overline{x'})=1,
\end{equation*}
one has
\steco
\begin{equation} \label{eq:XII.21} {} \prof F'_{x'}\geq n+1.
\end{equation}
Now this condition is satisfied at the points $x'$ where $F'$ is not free, since for such an $x'$ one has $x'\notin U_0'$ by virtue of d), hence $g$ is an isomorphism there, and by virtue of d) again, $F$ is of depth $\geq n+1$ at $g(x')$, hence $F'$ is of depth $\geq n+1$ at $x'$. It thus suffices to verify condition (\Ref{eq:XII.21}) at those $x'\in U'$ satisfying (\Ref{eq:XII.21}), and at which $F'$ is free. Now it suffices for this to prove that one has
\begin{equation*} \label{eq:XII.21bis} \tag{21 bis}
\prof \Oo_{X', x'}\geq n+1
\end{equation*}
%
at these points, \afortiori it suffices to establish that one has this relation at all the points~$x$ of $U'$ satisfying (\Ref{eq:XII.20} bis). Now, again by virtue of the criterion~\Ref{VIII.2.3} of Expos\'e \Ref{VIII}, this is equivalent to the assertion that the Modules
\[
\SheafH^p_{Y'}(\Oo_{X'}) \hspace{5mm} \text{for} \; p\leq n+1
\]
are coherent. In fact, we are going to prove that they are even zero, or what amounts to the same by virtue of Expos\'e~\Ref{III}, that one has
\steco
\begin{equation} \label{eq:XII.22} {} \prof \Oo_{X', x'}\geq n+2 \hspace{5mm} \text{for every }\; x'\in Y'.
\end{equation}
For this, we distinguish two cases. If $x'\notin X'_0$, then $f$ is an immersion at $x'$, and one must verify that $F$ is of depth $\geq n+2$ at the image $x=f(x')$, which is nothing other than condition c). If on the other hand $x'\in X'_0$, \ie $x=f(x')\in X_0$ hence $x\in Y\cap X_0$, one applies conditions a) and b) thanks to the

\begin{lemma} \label{XII.4.2}
Let $X$ be a locally noetherian prescheme, $X_0$ a closed subprescheme of $X$ defined by an ideal $\Jj$, $X'=\Spec(\bigoplus_{m\geq 0} \Jj^m)$, $X_0'=\Spec(\bigoplus_{m\geq 0} \Jj^m/\Jj^{m+1}) = X'\times_XX_0$, $x$ a point of $X_0$ at which $X$ is normally flat along $X_0$ \ie such that $\gr_{\Jj}(\OX)= \bigoplus_{m\geq 0}\Jj^m/\Jj^{m+1}$ is flat there as a Module on $X_0$. Then for every sequence of elements $f_i$ ($1\leq i\leq m$) of $\Oo_{X, x}$ whose images in $\Oo_{X_0, x}$ form an $\Oo_{X_0, x}$-regular sequence, and every $x'\in X'$ above $x$, the images of the $f_i$ in $\Oo_{X', x'}$ (\resp in $\Oo_{X'_0, x'}$) likewise form an $\Oo_{X', x'}$-regular sequence (\resp an $\Oo_{X'_0, x'}$-regular sequence), in particular one has
\steco
\begin{equation} \label{eq:XII.23}
\prof \Oo_{X', x'} \geq \prof \Oo_{X_0, x}, \quad \prof \Oo_{X'_0, x'} \geq \prof \Oo_{X_0, x}.
\end{equation}
\end{lemma}

To prove it, one may suppose that $X$ is local with closed point $x$, hence affine with ring $A=\Oo_{X, x}$, $\Jj$ being defined by the Ideal $J$, and it suffices to prove that for every sequence $f_i$ ($1\leq i\leq m$) of elements of $A$ whose images in $A/J$ form an $A/J$-regular sequence, the $f_i$ likewise form a $(\bigoplus_{m\geq 0} J^m)$-regular sequence and a $(\bigoplus_{m\geq 0} J^m/J^{m+1})$-regular sequence, \ie for every $m$, it forms a $J^m$-regular sequence and a $J^m/J^{m+1}$-regular sequence. The second assertion is trivial, since $J^m/J^{m+1}$ is a free module over $A/J$. The first follows, by looking at the $J$-adic filtration of $J^m$, and noting that for the graded module associated to $J^m$ for the said filtration, the sequence of the~$f_i$ is regular.

This proves \Ref{XII.4.2} and consequently \Ref{XII.4.1}~1).

Let us prove \Ref{XII.4.1} 2). For this, let us use the cartesian square
\[
\xymatrix{
X'_0\ar[d]_{f_0}&\ar[l]_-{i'_0} U'_0\ar[d]^{g_0} \\
X_0&\ar[l]_-{i_0} U_0
}
\]
and proceeding as in the beginning of the proof of 1), one finds that one has
\steco
\begin{equation} \label{eq:XII.24}
\tbigoplus_{m\geq 0} \R^pi_\ast(\Jj^m\formelF/\Jj^{m+1}\formelF) \simeq {f_0}_\ast (\R^p{i'_0}_\ast(\formelF_0')),
\end{equation}
where $\formelF_0=\formelF/\Jj\formelF$ and where $\formelF_0'=g_0^\ast(\formelF_0)$, (using the fact that $\formelF$ is locally free). Hence the conclusion of 2) is equivalent to saying that for $p\leq n$, $\R^p{i'_0}_\ast(\formelF_0')$ is a coherent Module.

%
Here again, taking into account that $\formelF_0'$ is locally free, the criterion \Ref{VIII}~\Ref{VIII.2.3} allows us to reduce to proving that this is so if one replaces $\formelF_0'$ by $\Oo_{U'_0}$, \ie to proving that the Modules
\[
\SheafH^p_{Y_0'}(\Oo_{X_0'}) \hspace{5mm} \text{for} \; p\leq n+1 \hspace{5mm} \text{(where} \; Y'_0=Y'\cap X_0'=X_0'- U'_0 \text{)}
\]
are coherent. One proves again that they are in fact zero, i.e.\ that one has
\steco
\begin{equation} \label{eq:XII.25} {}\prof \Oo_{X'_0, x'} \geq n+2 \hspace{5mm} \text{for every} \; x'\in Y'_0.
\end{equation}
Now this indeed follows from conditions a) and b), taking \Ref{XII.4.2} into account. This completes the proof of \Ref{XII.4.1}.
\skipqed
\end{proof}

\begin{remark} \label{XII.4.3}
One sees at once, by descent, that the hypothesis: $X$ locally immersible in a regular scheme, may be replaced by the following weaker one: there exists a morphism $\overline{X}\to X$, faithfully flat and quasi-compact, such that $\overline{X}$ is locally immersible in a regular scheme.
\end{remark}

Theorem \Ref{XII.4.1} puts us in a position to apply the results of Expos\'e \Ref{IX} (comparison and existence theorems). We shall be interested in particular in

\begin{corollary} \label{XII.4.4}
Suppose conditions a), b), c) of Theorem \Ref{XII.4.1} are satisfied, with $n=1$, and $X=\Spec(A)$, $A$ being separated and complete for the $J$-adic topology. Then
\begin{enumerate}
\item The functor $F\mto \widehat{F}$ from the category of coherent locally free Modules on~$U$ to the category of coherent locally free Modules on $\widehat{U}$ is fully faithful. \item For every coherent locally free Module $\formelF$ on $\widehat{U}$, there exists a coherent Module $F$ on $U$ and an isomorphism $\widehat{F}\isomto \formelF$.
\end{enumerate}
In particular, if for every $x\in U$ whose closure in $U$ does not meet $U_0$ \ie such that $\overline{x}\cap X_0\subset Y$, one has $\prof \Oo_{U, x}\geq 2$, then the pair $(U, U_0)$ satisfies the effective Lefschetz condition ($\Leff$) of Expos\'e \Ref{X}.
\end{corollary}
\noindent
(For the last assertion, one proceeds as in \Ref{X}~\Ref{X.2.1})

Special case of \Ref{XII.4.4}:

\begin{corollary} \label{XII.4.5}
Let $A$ be a noetherian ring, $J$ an ideal of $A$ contained in the radical, $A_0=A/J$. Assume
\begin{enumeratei}
\item
$\prof A_0\geq 3$.
\item
$\gr_J(A)$ is a free $A_0$-module.
\item
$A$ is complete for the $J$-adic topology.
\end{enumeratei}

Let $X=\Spec(A)$, $X_0=\Spec(A_0)=V(J)$, $a$ the closed point of $X$, $U=X-\{a\}$, $U_0=X_0-\{a\}$, $\widehat{U}$ the formal completion of $U$ along $U_0$. Then the functor $F\mto \widehat{F}$ from the category of coherent locally free Modules on $U$ to the category of coherent locally free Modules on $\widehat{U}$ is fully faithful. Moreover, for every coherent locally free Module $\formelF$ on $\widehat{U}$, there exists a coherent Module (not necessarily locally free!) $F$ on $U$, and an isomorphism $\widehat{F}\simeq\formelF$.
\end{corollary}

One will note that thanks to \Ref{XII.4.3}, we have not had to assume that $A$ is a quotient of a regular ring, for the completion of $A$ for the $\rr(A)$-adic topology satisfies this condition in any case.

Proceeding as in Expos\'es \Ref{X} and \Ref{XI}, one concludes from \Ref{XII.4.5}

\begin{corollary} \label{XII.4.6}
Under the conditions of \Ref{XII.4.5}, one has the following:
\begin{enumeratea}
\item
$U$ and $U_0$ are connected (\Ref{III}~\Ref{III.3.1}).

Choosing a geometric base-point in $U_0$, the homomorphism
\[
\pi_1(U_0)\to \pi_1(U)
\]
is surjective.
\item
The homomorphism
\[
\Pic(U)\to \Pic(U_0)
\]
is injective.
\end{enumeratea}
\end{corollary}

To prove b), taking \Ref{XII.4.5} into account, this amounts to verifying that every isomorphism $L_0'\isomto L_0$ lifts to an isomorphism $\widehat{L'}\isomto \widehat{L}$. Now for this one lifts step by step to isomorphisms $L_n'\isomto L_n$, the obstructions lying in $\H^1(U_0, \Jj^n/\Jj^{n+1})$, and these modules are zero from the fact that $J^n/J^{n+1}$ is free and $\prof A_0\geq 3$.

We are now in a position to prove the
%
\begin{theorem} \label{XII.4.7}
Let $A$ be a noetherian local ring, $J$ an ideal of $A$ contained in its radical, $A_0=A/J$. Assume
\begin{enumeratei}
\item
$\prof A_0\geq 3$.
\item
$\gr_J(A)$ is a free module over $A_0$.
\end{enumeratei}
\noindent
Then, if $A_0$ is ``pure'' (\Ref{X}~\Ref{X.3.1}) (\resp parafactorial (\Ref{XI}~\Ref{XI.3.1})), the same is true of~$A$.
\end{theorem}

\begin{proof}
By descent, one may suppose that one also has
\begin{enumeratei}\setcounter{enumi}{2}
\item
$A$ is complete for the $J$-adic topology.
\end{enumeratei}

Indeed, by virtue of (i) and (ii), one has $\prof(A)\geq 3$ hence $\prof(\widehat{A})\geq 3$, where $\widehat{A}$ is the completion of $A$ for the $J$-adic topology, and one applies \Ref{X}~\Ref{X.3.6} and \Ref{XI}~\Ref{XI.3.6}. One is thus under the conditions of \Ref{XII.4.5}. As $\prof(A)\geq 3\geq 2$, to say that $A$ is parafactorial simply means that $\Pic(U)=0$, and by virtue of \Ref{XII.4.6} b) it suffices for this that $\Pic(U_0)=0$, \ie that $A_0$ be parafactorial. To prove that $A$ is ``pure'' if $A_0$ is, one must prove that if $V$ is an \'etale covering of $U$, defined by an algebra $\Bb$ on~$U$, then $\H^0(U, \Bb)$ is a finite \'etale algebra over $A$. Now $A_0$ being pure, the same is true of the $A_n$ (which differ from it only by nilpotent elements), hence for every $n$, $B_n=\H^0(U, \Bb_n)$ is an \'etale algebra over $A/J^{n+1}$, and of course these algebras glue, so that $\varprojlim B_n$ is an \'etale algebra over $A$. Now by virtue of \Ref{XII.4.5}, this algebra is none other than $\H^0(U, \Bb)$, which establishes our assertion.
\skipqed
\end{proof}

\begin{corollary} \label{XII.4.8}
Let $f\colon X\to Y$ be a flat morphism of locally noetherian preschemes, $x\in X$, $y=f(x)$, suppose that $\Oo_{X_y, x}$ is a ``pure'' (\resp parafactorial) local ring of depth $\geq 3$, then the same is true of $\Oo_{X, x}$.
\end{corollary}

This is the result of the type promised at the end of the preceding \numero, in order to generalize Corollaries \Ref{XII.3.5} and those following. One thus finds, using \Ref{XII.3.4}, the

\begin{corollary} \label{XII.4.9}
Let $f\colon X\to S$ be a projective and flat morphism, with $S$ locally noetherian, $\OX(1)$ an invertible Module on $X$ ample with respect to $S$, $t$ a section of $\OX(1)$ such that for every $s\in S$, the section $t_s$ induced on $X_s$ is $\Oo_{X_s}$-regular, $X_0$ the subscheme of zeros of $t$, $X_m$ the subscheme of zeros of $t^{m+1}$. Assume that for every $s\in S$, $X_s$ is of depth $\geq 3$ at all its closed points. Then:
\begin{enumeratea}
\item
If the local rings of the closed points of $X_s- X_{0, s}$ ($s\in S$) are ``pure'', for example are complete intersections, then the functor $X'\mto X_0'=X'\times_X X_0$ from the category of \'etale coverings of $X$ to the category of \'etale coverings of $X_0$ is an equivalence of categories; in particular, choosing a geometric base-point in $X_0$, the homomorphism
\[
\pi_1(X_0)\to \pi_1(X)
\]
is an isomorphism\nde{\label{FH}let us point out
the spectacular connectedness result obtained since by Fulton and
Hansen, in the case where $S=\Spec(k)$ ($k$ an algebraically
closed field). Let $g\colon X\to {\PP^m_k}\times{\PP^m_k}$ be such that $\dim
g(X)>m$; then the inverse image of the diagonal is connected. This
allows one among other things to generalize
Corollary~\Ref{XII.4.9} when $f$ is the structural morphism
of the projective space $\PP^m_k$ over $S=\Spec(k)$: precisely, an
irreducible subvariety $X$ of ${\PP^m_k}$ of dimension $>m/2$ has
trivial fundamental group! (\cf Fulton~W.\ \& Hansen~J., ``
A connectedness theorem for projective varieties, with
applications to intersections and singularities of mappings'',
\emph{Ann. of Math. (2)} \textbf{110} (1979), \numero 1, p\ptbl
159--166). For generalizations to the case of Grassmannians or
abelian varieties, see {Debarre~O.}, ``Th\'eor\`emes de
connexit\'e pour les produits d'espaces projectifs et les
grassmanniennes'', \emph{Amer. J.~Math.} \textbf{118} (1996),
\numero 6, p\ptbl 1347--1367 and ``Th\'eor\`emes de connexit\'e et
vari\'et\'es ab\'eliennes'', \emph{Amer. J.~Math.} \textbf{117}
(1995), \numero 3, p\ptbl 787--805. The result on the triviality of the
fundamental group of $X$ as above was obtained
independently by Faltings, who proves moreover that
$\text{Pic}(X)$ has no torsion prime to the characteristic
of $k$, this by methods of algebraization of formal bundles,
more in the line of Grothendieck's techniques, \cf
(Faltings~G., ``Algebraization of some formal vector bundles''
\emph{Ann. of Math. (2)} \textbf{110} (1979), \numero 3, p\ptbl
501--514).}

\item
If the local rings of the closed points of the $X_s- X_{0, s}$ ($s\in S$) are ``parafactorial'', for example regular, or complete intersections of dimension $\geq 4$, then for every integer $m$ such that $\R^if_{0\ast}(\Oo_{X_0}(-n))=0$ for $n>m$ and $i=1, 2$, the map $\Pic(X)\to \Pic(X_m)$ is bijective.

Moreover if $S$ is noetherian, and the $X_{0, s}$ are of depth $\geq 3$ at their closed points, there exist such $m$ (\cf \Ref{XII.1.5}).
\end{enumeratea}
\end{corollary}

\begin{remark} \label{XII.4.10}
Under the conditions of the last assertion of \Ref{XII.4.9} b), one has seen in \Ref{XII.1.5} that there exists an $m$ such that $n>m$ even implies $\H^i(X_0, \Oo_{X_0}(-n))=0$ for $i=1, 2$, and even for $i\leq 2$). This condition is stronger than $\R^if_{\ast}(\Oo_{X_0}(-n))=0$ for $i=1, 2$, and it has moreover the advantage of being stable under change of base. The same is true of the depth hypotheses made in \Ref{XII.4.9}, and likewise of a hypothesis of the type ``the $X_s$ are locally complete intersections''. It then follows, under these conditions that \Ref{XII.4.9} b) also implies that the morphism of functors
\[
\SheafPic_{X/S}\to \SheafPic_{X_n/S}
\]
in $\catSch_{/S}$ is an isomorphism, hence also the morphism for the relative Picard schemes, when these exist:
\[
\Pic_{X/S}\to \Pic_{X_n/S}.
\]
Even in the case where $S$ is the spectrum of an algebraically closed field, this statement is distinctly more precise than the statement consisting in saying only that $\Pic(X)\to \Pic(X_n)$ is bijective.

One may ask whether one can always take $n=0$ in the preceding conclusions (thus assuming the $X_{0, s}$ are of depth $\geq 3$ at their closed points). When $X_0$ is smooth over $S$ and the residual characteristics of $S$ are zero, this is indeed the case, by virtue of Kodaira's ``vanishing theorem'', (proved by transcendental means, using a K\"ahler metric) which implies that for every connected smooth projective scheme of dimension $n$ over a field $k$ of characteristic zero, and every ample invertible Module $L$ on $X$, one has $\H^i(X, L^{-1})=0$ for $i\neq n$. It is not known\nde{as Raynaud has remarked, the decomposition result for the de Rham complex of Deligne and Illusie easily implies the vanishing of the group $H^i(X, L^{-1})$ (with $L$ ample on $X$ projective smooth over $k$ of characteristic $p>0$) for $i<\inf(p, \dim(X))$ as soon as one assumes that $X$ is liftable to a flat scheme over $W_2(k)$ (\cf Deligne~P.\ \& Illusie~L., ``Rel\`evements modulo~$p^2$ et d\'ecomposition du complexe de de Rham'', \emph{Invent. Math.} \textbf{89} (1987), \numero 2, p\ptbl 247--270); this gives a purely algebraic proof of Kodaira's result for projective varieties in characteristic zero. If $X$ is not liftable, it is well known that the ``vanishing theorem'' is false; \cf the example in (Raynaud~M., ``Contre-exemple au "vanishing theorem" en caract\'eristique $p>0$'', in \emph{C.P.~Ramanujam---a tribute}, Tata Inst. Fund. Res. Studies in Math., vol.~8, Springer, Berlin-New York, 1978, p\ptbl 273--278); see also the very pretty examples in (Haboush~W.\ \& Lauritzen~N., ``Varieties of unseparated flags'', in \emph{Linear algebraic groups and their representations (Los Angeles, CA, 1992)}, Contemp. Math., vol.~153, American Mathematical Society, Providence, RI, 1993, p\ptbl 35--57), simplified in (Lauritzen~N.\ \& Rao~A.P., ``Elementary counterexamples to Kodaira vanishing in prime characteristic'', \emph{Proc. Indian Acad. Sci. Math. Sci.} \textbf{107} (1997), \numero 1, p\ptbl 21--25). On the other hand, I do not know an example where the arrow $\Pic(X_{n+1})\to\Pic(X_n)$ is not surjective for $n>1$ in positive characteristic, where $X_n$ denotes a thickened hyperplane section of $X$ projective smooth as above.} at the present time whether this theorem can be replaced by a generalization in characteristic $p>0$, and whether the smoothness hypothesis can be replaced by a hypothesis of a more general nature (bearing on the depth, or of ``complete intersection'' type...).
\end{remark}

% English body fragment for SGA 2 XII (en-4.tex)

\section[Universal finiteness conditions]
{Universal finiteness conditions for a non-proper morphism} \label{XII.5} Let us recall for the record the \setcounter{equation}{25}

\begin{proposition} \label{XII.5.1}
Let $f:X\to S$ be a proper morphism of preschemes, with $S$ locally noetherian, $U$ an open subset of $X$, $g:U\to X$ the canonical immersion, $h=fg:U\to S$, $F$ a Module on $U$. Suppose that the Modules $\R^ig_\ast(F)$ are coherent for $i\leq n$ (hypothesis of a local nature on $X$, which is verified in practice with the help of the criterion \Ref{VIII}~\Ref{VII.2.3}). Then $\R^ih_\ast(F)$ is coherent for $i\leq n$.
\end{proposition}

This follows at once from the Leray spectral sequence
\steco
\begin{equation} \label{eq:XII.26} {} \E_2^{p, q}=\R^pf_\ast(\R^qg_\ast(F)) \To \R^\ast h_\ast(F),
\end{equation}
and from the fact that the higher direct images under $f$ of a coherent Module on $X$ are coherent (\EGA III 3.2.1).

\begin{proposition} \label{XII.5.2}
Let $S$ be a locally noetherian prescheme, $\Ss$ a quasi-coherent graded Algebra, of finite type over $S$, generated by $\Ss_1$, $X$ a subprescheme of $\Proj(\Ss)$, $\OX(1)$ the invertible Module on $X$ very ample relative to $S$ induced by $\Proj(\Ss(1))$, $U$ an open subset of $X$, $g:U\to X$ the canonical immersion, $h=fg:U\to S$, $F$ a quasi-coherent Module on $U$, whence twisted Modules $F(m)=F\otimes \OX(m)$) ($m\in\ZZ$), $n$ an integer, $m_0$ an integer. The following conditions are equivalent:
\begin{enumeratei}
\item
$\R^ig_\ast(F)$ is coherent for $i\leq n$.
\item
$\bigoplus_{m>m_0}\R^ih_\ast(F(m))$ is an $\Ss$-Module of finite type for $i\leq n$.
\end{enumeratei}
\end{proposition}

\begin{proof}
Replacing in the spectral sequence above $F$ by $F(m)$ one
finds a spectral sequence of graded $\Ss$-Modules
\[
\E_2^{p, q}= \tbigoplus_{m\geq m_0} \R^pf_\ast(\R^qg_\ast(F(m)) \To \tbigoplus_{m\geq m_0} \R^\ast h_\ast(F(m)).
\]
As one has
\[
\R^qg_\ast(F(m)) \simeq \R^qg_\ast(F) (m),
\]
one sees that if the $\R^ig_\ast(F)$ are coherent, $\E_2^{p, q}$ is of finite type over $\Ss$ for $q\leq n$, thanks to part a) of lemma \Ref{XII.5.3} below, which implies that the abutment is of finite type over $\Ss$ in degree $i\leq n$. This proves (i) \ALORS (ii). Moreover, reasoning in the abelian category of graded $\Ss$-Modules modulo the thick subcategory $\mathcal{C}$ of those which are quasi-coherent of finite type, one finds by the preceding spectral sequence
\[
\tbigoplus_{m\geq m_0} \R^{n+1} h_\ast(F(m)) \simeq \tbigoplus_{m\geq m_0} f_\ast (\R^{n+1} g_\ast(F)(m)) \mod \mathcal{C},
\]
which proves that if the left-hand side is an $\Ss$-Module of finite type, then $\R^{n+1} g_\ast(F)$ is coherent, by virtue of part b) of lemma \Ref{XII.5.3}. This proves the implication (ii) \ALORS (i) by induction on $n$. It remains to prove:

\begin{lemma} \label{XII.5.3} Let $S$, $\Ss$, $X$, $f$ be as in \Ref{XII.5.2}, and $G$ a quasi-coherent Module on $X$, $m_0$ an integer. Then
\begin{enumeratea}
\item
If $G$ is coherent, then for every integer $i$, the graded Module
\[
\tbigoplus_{m\geq m_0} \R^i f_\ast(G(m))
\]
over $\Ss$ is of finite type.
\item
Conversely, suppose that the Module $\bigoplus_{m\geq m_0} \R^i f_\ast(G(m))$ over $\Ss$ is of finite type, then $G$ is coherent.
\end{enumeratea}
\end{lemma}

\begin{proof}[Proof of \Ref{XII.5.3}]
For a), the case $i=0$ is given in \EGA III 2.3.2, the case $i>0$ in \EGA III 2.2.1 (i)(ii) which says that the $ \R^i f_\ast G(m)$ are coherent, and vanish for $m$ large (if one assumes $S$ noetherian, which is permissible). For b), one notes that $G$ is isomorphic to $\SheafProj (\bigoplus_{m\geq m_0}f_\ast(G(m))$ (\EGA II 3.4.4 and 3.4.2), which proves that $G$ is coherent if $\bigoplus_{m\geq m_0}f_\ast(G(m))$ is of finite type over $\Ss$, by virtue of \loccit 3.4.4.
\skipqed
\end{proof}
\skipqed
\end{proof}

\begin{corollary}[{{\normalfont$S$ noetherian}}] \label{XII.5.4}
Suppose that $\R^ig_\ast(F)$ is coherent for $i\leq n$, then for $i\leq n+1$, and $m$ large, one~has a canonical isomorphism:
\[
\R^ih_\ast(F(m)) \simeq f_\ast(\R^ig_\ast(G)(m)).
\]
\end{corollary}

Indeed the spectral sequence (\Ref{eq:XII.26}) for $F(m)$ then degenerates in degree $\leq n$, by \EGA III 2.2.1~(ii), whence at once the result (which moreover recovers the implication (ii) \ALORS (i) of~5.2).

%
\begin{corollary} \label{XII.5.5}
Under the preliminary conditions of \Ref{XII.5.2}, $S$ noetherian, the following conditions are equivalent:
\begin{enumeratei}
\item
$\bigoplus_{m\geq m_0}h_\ast(F(m)$ is of finite type over $S$, and $\R^ih_\ast(F(m))=0$ for $0<i\leq n$ and $m$ large.
\item
$g_\ast(F)$ is coherent, and $\R^ig_\ast(F)=0$ for $0<i\leq n$.
\item[\textup{(ii bis)}] $g_\ast(F)$ is coherent, and $\prof_Y g_\ast(F)>n+1$.
\end{enumeratei}
\end{corollary}

The equivalence of (ii) and (ii bis) is contained in \Ref{III}~\Ref{III.3.3}. Moreover, by virtue of \Ref{XII.5.2} the conditions (i) and (ii) both imply that the $\R^ig_\ast(F)$ ($i\leq n$) are coherent. The equivalence of (i) and (ii) then follows from \Ref{XII.5.4}, taking into account the fact that every a coherent Module $G$ on $X$, one~has $G=0$ if and only if $f_\ast(G(m))=0$ for $m$ large, for example by virtue of \EGA III 2.2.1 (iii).

\begin{remark} \label{XII.5.6}
One can interpret the criteria \Ref{XII.5.2} and \Ref{XII.5.5} by saying that the \og simultaneous finiteness condition\fg \Ref{XII.5.2} (ii) is expressed by properties of local regularity (in terms of depth thanks to \Ref{VIII}~\Ref{VIII.2.1}) of $F$ at the points of $U$ neighbouring $Y=X- U$, whereas the \og asymptotic vanishing condition\fg \Ref{XII.5.5} (i) is of a distinctly stronger nature, and is expressed by conditions of local regularity of $g_\ast(F)$ at the points of~$Y$ itself. It would be interesting, in order to generalize the Lefschetz-type theorems for projective morphisms to quasi-projective morphisms, to find local criteria on $X$ necessary and sufficient for the $\Ss$-Modules $\bigoplus_{m\geq 0} \R^ih_\ast(F(m))$ for $i\leq n$ to be of finite type. When $S$ is the spectrum of a field (and no doubt more generally, if it is the spectrum of an artinian ring) and $Y=X- U$ is finite, one can show that it is necessary and sufficient that the following conditions be satisfied:
\begin{enumerate}
\item[\textup{1\noo)}] $\prof F_x >n$ for every closed point $x$ of $U$ (compare \Ref{XII.1.4}).
\item[\textup{2\noo)}] $\R^ig_\ast(F)$ is coherent for $i\leq n$, or what amounts to the same, there exists an open neighbourhood $V$ of $Y$ such that for every closed point $x$ of $U\cap V$, one has $\prof F_x >n+1$.
\end{enumerate}
\end{remark}

\begin{proposition} \label{XII.5.7}
Let $S$ be a locally noetherian prescheme, $g:U\to X$ a morphism of preschemes of finite type over $S$\sfootnote{It in fact suffices that $g$ be quasi-compact and quasi-separated (\textup{\EGA IV1.2.1}), with no condition on $U, X$.}, with structural morphisms $h$ and $f$, $F$ a quasi-coherent Module on $U$, $n$ an integer. The following conditions are equivalent:
\begin{enumeratei}
\item
For every base change $S'\to S$, with $S'$ noetherian, the module $\R^ng'_\ast(F')$ on $X'$ is coherent.
\item
For every base change as above, and every coherent Ideal $\Jj$ on~$S'$, denoting by $\Ii$ the Ideal $\Jj\Oo_{X'}$ on $X'$, the graded Module
\[
\tbigoplus_{m\geq 0} \R^ng'_\ast(\Ii^mF')
\]
over $\bigoplus_{m\geq 0} \Ii^m$ is of finite type.
\item
For every base change $S'\to S$, and $\Jj$ as above, the graded Module
\[
\tbigoplus_{m\geq 0} \R^ng'_\ast(\Ii^mF'/\Ii^{m+1}F')
\]
over $\gr_I(\Oo_{X'})=\bigoplus_{m\geq 0} \Ii^m/\Ii^{m+1}$ is of finite type.
\end{enumeratei}
\end{proposition}

Obviously (ii) \ALORS (i) and (iii) \ALORS (i), as one sees by setting $\Ii=0$ in conditions (ii) and (iii). The converse implications are obtained by applying (i) to the composite base change $S^{''}\to S'\to S$, where $S^{''}$ is equal to $\Spec \bigoplus_{m\geq 0} \Jj^m$ \resp $\Spec \bigoplus_{m\geq 0} \Jj^m/\Jj^{m+1}$.

The interest of this proposition is that the conditions of the form (ii) are those which intervene in the \og algebraic-formal comparison theorems\fg, whereas the conditions of the form (iii) intervene in the \og existence theorems\fg which complete them, \cf expos\'e \Ref{IX}. A first interesting case is that where $f:X\to S$ is the identity, and where it is thus a matter of conditions on a morphism $h:U\to S$ locally of finite type and a Module $F$ quasi-coherent on $U$ flat with respect to $S$. To obtain sufficient conditions, we shall assume that $U$ embeds via $g:U\to X$, as an open subprescheme of an $X$ \emph{proper} over $S$. Applying \Ref{XII.5.1}, one thus sees:

\begin{corollary} \label{XII.5.8}
Let $f:X\to S$ be a proper morphism, with $S$ locally noetherian, $U$ an open set of $X$, $g:U\to X$ the canonical immersion, $h=fg:U\to S$, $F$ a quasi-coherent Module on $X$, flat with respect to $S$. Suppose that for every base change $S'\to S$, with $S'$ locally noetherian, one has $\R^ig'_\ast(F')$ coherent on $X'$ for $i\leq n$. Then one~has the following:
\begin{enumeratei}
\item
For every base change $S'\to S$, with $S'$ locally noetherian, $\R^ih'_\ast(F')$ is coherent on $S'$ for $i\leq n$.
\item
For every $S'\to S$ as above, and every coherent Ideal $\Jj$ on $S'$, the graded Modules
\[
\tbigoplus_{m\geq 0} \R^ih'_\ast(\Jj^mF')
\]
over $\bigoplus_{m\geq 0}\Jj^m$ are of finite type for $i\leq n$.
\item
For every $S'\to S$ and $\Jj$ as above, the graded Modules
\[
\tbigoplus_{m\geq 0} \R^ih'_\ast(\Jj^mF'/\Jj^{m+1}F')
\]
over $\gr_\Jj(\Oo_{S'})=\bigoplus_{m\geq 0}\Jj^m/\Jj^{m+1}$ are of finite type for $i\leq n$.

Moreover, under the conditions of (ii), and by virtue of the comparison theorem \Ref{IX.1.1}, denoting by $\widehat{S'}$ the formal completion of $S'$ with respect to $\Jj$, and by $\widehat{U'}$ that of $U'$ with respect to $\Jj\Oo_{U'}$, the canonical homomorphisms
\[
\widehat{\R^ih'_\ast(F')} \to \R^i\widehat{h'}_\ast(\widehat{F'}) \to \varprojlim_k\R^ih'_\ast(F'_k)
\]
are isomorphisms for $i\leq n-1$.
\end{enumeratei}
\end{corollary}

\begin{remark} \label{XII.5.9}
Suppose that one is moreover under the conditions of \Ref{XII.5.8} with $F$ \emph{coherent}, and consider a base change $S'\to S$ as in \Ref{XII.5.9} (i). Suppose moreover that $S'$ is locally immersible in a regular scheme, or more generally, that there exists a morphism $S^{''}\to S$ faithfully flat and quasi-compact, such that $S^{''}$ is locally immersible in a regular scheme; this condition is satisfied in particular if $S'$ is local. Then the conclusion of \Ref{XII.5.8} (i) and (ii) remains valid when one replaces $F'$ therein by a Module $G'$ on $U'$, such that every point of $U'$ has an open neighbourhood on which $G'$ is isomorphic to a Module of the form ${F'}^n$. Indeed, one is reduced to the case $S'$ itself is locally immersible in a regular scheme, so that the same holds for $S^{''}=\Spec\big(\bigoplus_{m\geq 0}\Jj^m\big)$ and for $X^{''}=X'\times_{S'} S^{''}= X\times_{S} S^{''}$, which are of finite type over it. One then applies the finiteness criterion \Ref{VIII}~\Ref{VIII.2.3} for the direct images for $i\leq n$ of $G^{''}$ under the immersion $U^{''}\to X^{''}$, noting that they are satisfied by hypotheses for $F^{''}$, hence also for $G^{''}$ since they are expressed in terms of depth and $G^{''}$ is locally isomorphic to a ${F^{''}}^n$. The same argument shows that if $\Gg'$ is a coherent Module on $\widehat{U'}$ (completion of $U'$ with respect to the Ideal $\Jj\Oo_{U'}$), such that $\Gg'_0=\Gg'/\Jj\Gg'$ is locally of the form $F_0^{'n}$, then the conclusion of (iii) remains valid upon replacing $F'$ therein by $\Gg'$. One thus obtains the following result, using the results of expos\'e \Ref{IX}:
\end{remark}

\begin{corollary} \label{XII.5.10}
Let $f:X\to S$ be a proper morphism, with $S$ locally noetherian, $U$ an open subset of $X$; one assumes $U$ flat with respect to $S$, and that for every base change $S'\to S$, with $S'$ locally noetherian, one has $R^ig'_\ast(\Oo_{U'})$ coherent on $X'$ for $i=0, 1$. Suppose then that $S'$ is of the form $\Spec(A)$, where $A$ is noetherian ring equipped with an ideal $J$ such that $A$ is separated and complete for the $J$-adic topology. Under these conditions:
\begin{enumeratei}
\item
The functor $F\mto\widehat{F}$ from the category of locally free Modules on $U'$ to the category of locally free Modules on $\widehat{U'}$ is fully faithful.
\item
For every locally free Module $\formelF$ on $\widehat{U'}$, there exists a coherent Module $F$ on~$U'$ (not necessarily locally free, alas), and an isomorphism $\widehat{F}\simeq \formelF$.
\end{enumeratei}
\end{corollary}

It remains only to prove (ii), thanks to \Ref{XII.5.9}. Now according to this remark and \Ref{IX.2.1} it follows that $\formelF$ is induced by a coherent Module $\formelG$ on $\widehat{X'}$. By the existence theorem \EGA III 5.1.4, $\formelG$ is of the form $\widehat{F}$, where $F$ is coherent on $X$, whence the conclusion.

\begin{remarks} \label{XII.5.11}
1\textsuperscript{o}) Using \Ref{XII.5.10}, \Ref{XII.4.7} and a suitable hypothesis, saying that certain local rings of the geometric fibres of $X'\to S'$ are \og pure\fg \resp parafactorial, one should be able to obtain statements saying that the functor $Z'\mto \widehat{Z'}$ from the category of \'etale coverings of $X'$ to the category of \'etale coverings of $\widehat{X'}$ (or what amounts to the same, of $X_0'$) is an equivalence of categories, \resp that the functor $L\mto \widehat{L}$ from the category of invertible Modules on $X'$ to the category of invertible Modules on $\widehat{X'}$ is an equivalence. Using recent results of Murre, it is probable that one should be able to deduce existence theorems for Picard schemes for certain \emph{non-proper} algebraic schemes\nde{of course, in the projective case one will refer to Grothendieck's existence theorems of FGA; \cf Grothendieck~A., {\og Technique de descente et th\'eor\`emes d'existence en g\'eom\'etrie alg\'ebrique. VI. Les sch\'emas de Picard: propri\'et\'es g\'en\'erales\fg}, in \emph{S\'eminaire Bourbaki}, vol.~7, Soci\'et\'e math\'ematique de France, Paris, 1995, \Exp 236, p\ptbl 221--243 and {\og Technique de descente et th\'eor\`emes d'existence en g\'eom\'etrie alg\'ebrique. V. Les sch\'emas de Picard: th\'eor\`emes d'existence\fg}, in \emph{S\'eminaire Bourbaki}, vol.~7, Soci\'et\'e math\'ematique de France, Paris, 1995, \Exp 232, 143--161. The nine finiteness conjectures which are found there are proved in expos\'es \Ref{XII} and \Ref{XIII} of Mme Raynaud and of Kleiman of \SGA 6. For an excellent elementary text on the subject, see Kleiman's survey article (Kleiman~S., {\og The Picard scheme\fg}, to appear in Contemp. math.). For an application of these techniques to the global generalized Jacobians of a relative smooth curve, see (Contou-Carr\`ere~C., {\og La jacobienne g\'en\'eralis\'ee d'une courbe relative; construction et propri\'et\'e universelle de factorisation\fg}, \textit{C.~R. Acad. Sci. Paris S\'er.~A-B} \textbf{289}, (1979), \numero3, A203--A206 and {\og Jacobiennes g\'en\'eralis\'ees globales relatives\fg}, in \emph{The Grothendieck Festschrift, Vol. II}, Progr. Math., vol.~87, Birkh\"auser, Boston, 1990, p\ptbl 69--109). See also by the same author, in the purely local setting, the construction and the study of the functor \og local generalized Jacobian\fg ({\og Jacobienne locale, groupe de bivecteurs de Witt universel, et symbole mod\'er\'e\fg}, \textit{C.~R. Acad. Sci. Paris S\'er.~I Math.}  \textbf{318} (1994), \numero 8, p\ptbl 743--746). Moreover, if in the case of a projective and smooth morphism, the connected components of the Picard scheme are proper, this is no longer so in the singular case. There naturally arises the problem of the compactification of Picard schemes: this problem has been studied in detail, notably in (Altman~A.B. \& Kleiman~S., {\og Compactifying the Picard scheme\fg}, \emph{Adv. in Math.} \textbf{35} (1980), \numero 1, p\ptbl 50--112. and {\og Compactifying the Picard scheme. II\fg}, \emph{Amer. J.~Math.} \textbf{101} (1979), \numero 1, p\ptbl 10--41). The case of curves had been studied previously (D'Souza~C., {\og Compactification of generalised Jacobians\fg}, \emph{Proc. Indian Acad. Sci. Sect. A Math. Sci.} \textbf{88} (1979), \numero 5, p\ptbl 419--457). One even knows exactly when the compactified Jacobian of a curve is irreducible (Rego~C.J., {\og The compactified Jacobian\fg} \emph{Ann. Sci. \'Ec. Norm. Sup. (4)} \textbf{13} (1980), \numero 2, p\ptbl 211--223, the latter being the closure of the (ordinary) Jacobian when the curve is geometrically integral and locally planar; for a construction in families of compactified Jacobians, see (Esteves~E., {\og Compactifying the relative Jacobian over families of reduced curves\fg}, \emph{Trans. Amer. Math. Soc.} \textbf{353} (2001), \numero 8, p\ptbl 3045--3095). Since then, the existence results for the Picard scheme in the proper case have progressed since the original edition of \SGA 2; \cf (Murre~J.P., {\og On contravariant functors from the category of pre-schemes over a field into the category of abelian groups (with an application to the Picard functor)\fg}, \emph{Publ. Math. Inst. Hautes \'Etudes Sci.} \textbf{23} (1964), p\ptbl 5-43) and above all (Artin~M., {\og Algebraization of formal moduli. I\fg}, in \emph{Global Analysis (Papers in Honor of K\ptbl Kodaira)}, Univ. Tokyo Press, Tokyo, 1969, p\ptbl 21--71). See also (Raynaud~M., {\og Sp\'ecialisation du foncteur de Picard\fg}, \emph{Publ. Math. Inst. Hautes \'Etudes Sci.} \textbf{38} (1970), p\ptbl 27--76) in the case of a proper scheme over a discrete valuation ring, but not necessarily flat. For a discussion of more recent results, in particular those of Artin, for the Picard functor of proper and flat schemes, in particular in the case cohomologically flat in dimension $0$, see Chapter VIII of (Bosch~S., L\"utkebohmert~W. \& Raynaud~M., \emph{N\'eron models}, Ergebnisse der Mathematik und ihrer Grenzgebiete (3), vol.~21, Springer-Verlag, Berlin, 1990) and the references cited. Much more recently, very fine results have been obtained in the case of relative curves $f:X\to S$ over the spectrum $S$ of a discrete valuation ring with perfect residue field. More precisely, one assumes that $f$ is proper and flat, $X$~regular and $f_*\OX=\Oo_S$. On the other hand, one does not assume $f$ cohomologically flat in dimension~$0$, \ie one does not assume $H^1(X,\Oo)$ torsion-free. The Picard scheme is then not representable, whether by a scheme or an algebraic space, as soon as the torsion in question is non-zero. Let~$J$ be the N\'eron model of the generic fibre of $f$: it is a quotient of the Picard functor $P$. Then, Raynaud has shown that the kernel of the tangent map $H^1(X,\Oo)=\mathrm{Lie}(P)\to \mathrm{Lie}(J)$ coincides with the torsion subgroup of the $H^1$ and that the cokernel has the same length (see theorem 3.1 of (Liu~Q., Lorenzini~D. \& Raynaud~M., {\og N\'eron models, Lie algebras, and reduction of curves of genus one\fg}, \emph{Invent. Math.} \textbf{157} (2004), p\ptbl 455--518). This result allows the aforementioned authors to study the link between the Birch--Swinnerton-Dyer and Artin--Tate conjectures (see th\ptbl 6.6 of {\loccit}). As regards the local Picard scheme, see Boutot's thesis, cited in the editor's note~\eqref{Boutot}~page~\pageref{Boutot}.}. In a general way, the elimination of purity hypotheses in various existence theorems, notably of representability of functors such as the Hilbert functors, or Picard functors etc., with the help of the techniques developed in this seminar, deserves a systematic study.

2\textsuperscript{o}) One may propose to give manageable necessary and sufficient conditions, in terms of depth, for the universal finiteness condition envisaged in \Ref{XII.5.10} to be satisfied. When $S$ is the spectrum of a field, it follows easily from \EGA III 1.4.15 that it is necessary and sufficient that the $\R^ig_\ast(F)$ ($i\leq n$) be coherent, which is well expressed in terms of depth thanks to \Ref{VIII}~\Ref{VIII.2.3}. In the general case, one will note however that it does not suffice to require that the preceding condition be satisfied for all the fibres $U_s\subset X_s$ ($s\in S$), even in the case where $n=0$. Take for example $X=S$, $S$ the spectrum of a discrete valuation ring, $U$ the open set reduced to the generic point, $F=\Oo_U$.

3\textsuperscript{o}) Here however is a \emph{sufficient} condition ensuring that one is under the conditions of the hypothesis of \Ref{XII.5.10}: It suffices that $f$ be flat, and that for every $s\in S$ and every $x\in Y_s=X_s- U_s$, one have
\[
\prof \Oo_{X_s, x}\geq n+2.
\]
Indeed, taking into account lemma \Ref{XII.2.5} (\cf relation (\Ref{eq:XII.16}) after \Ref{XII.2.5}), it follows that one then has $g_\ast(\Oo_U)\simeq\OX$ and $\R^ig_\ast(\Oo_U)=0$ for $0<i\leq n$, and the same relations will obviously be satisfied after every base change $S'\to S$.
\end{remarks}

\end{document}
